\documentclass[10pt]{article}
\usepackage[margin=20mm]{geometry}
\usepackage{amsmath,amssymb,mathtools,fontspec,unicode-math,xurl,hyperref}
\IfFontExistsTF{DejaVu Serif}{\setmainfont{DejaVu Serif}}{\setmainfont{Latin Modern Roman}}
\setmathfont{latinmodern-math.otf}
\setlength{\parindent}{0pt}\setlength{\parskip}{6pt}\sloppy
\title{Technical Review of
The Density Finite Sums Theorem}
\author{AI-assisted technical review}
\date{}
\begin{document}
\maketitle
Abstract: This review concerns the exact acquired 31-page publisher PDF of Kra, Moreira, Richter, and Robertson. It identifies a false density assertion in the motivating weak-mixing example and two missing-hypothesis problems in the stated uniformity estimates. Explicit counterexamples establish the failures. The required restricted uses are examined separately: the weak-mixing example retains progressiveness, normalized estimates retain the zero-seminorm argument, and the actual multiple-average coefficients are distinct. The report also records positive checks of the correspondence, progressive-measure construction, recurrence reduction, and final pronilfactor argument. No counterexample to the density finite sums theorem is established. Scientific adjudication remains pending.\par
\section*{Identity, version, and review boundary}
The paper is The density finite sums theorem, by Bryna Kra, Joel Moreira, Florian K. Richter, and Donald Robertson, DOI https://doi.org/10.1007/s00222-025-01371-8 [paper]. Every source location below refers to the acquired publisher PDF, whose printed and PDF page numbers agree. Its first page reads Inventiones mathematicae (2025), 243:1–31, and records online publication on September 24, 2025. The publisher’s current landing page instead assigns the issue to January 2026 [record]. This bibliographic year difference is recorded without changing the supplied 2025 assignment or asserting equivalence to a newly downloaded file. The exact source hash, rather than a year label, identifies the reviewed object.\par
The preserved index contains only four target source/identity/manifest entries and no recoverable substantive report for this paper. The present work is therefore its necessary initial review. It is not an independent acceptance review. All 31 page texts were read in the root scientific conversation; bounded native helpers additionally checked the structural, uniformity, and final-progression arguments. Actual source images were used to resolve formulas and extraction artifacts. The scope is broader than an abstract-only assessment, but it is not an independent reconstruction of the complete Host–Kra structure theory or of every cited theorem.\par
\section*{Quantifiers and the proof architecture}
Theorem 1.1, p. 2, fixes an arbitrary finite \(k\). It produces a strictly increasing positive sequence \(b\) and a shift \(t\geq0\) such that every sum of \(m\) distinct selected terms, \(1\leq m\leq k\), lies in \(A-t\). Neither a single sequence working for all \(k\) nor an unshifted conclusion is claimed. For example, odd integers already rule out the unshifted conclusion when \(k\geq2\). No repeated-index sum is part of the assertion.\par
The dynamical statement requires full measure on the forward orbit closure of \(a\). Its reduction uses a Følner sequence in \(\mathbb N\) along which \(a\) is generic. This does not mean that \(a\) must be generic along the initial intervals, or that every open-set return sequence has density one. These distinctions matter in Problem 1.\par
The essential chain is Theorem 1.1 \(\leftarrow\) Theorem 1.2 \(\leftarrow\) Theorem 2.4 \(\leftarrow\) Theorem 3.3. Proposition 3.2 extracts an Erdős progression from a progressive measure concentrated on \(\{a\}\times X^k\). Theorems 4.8 and 4.9 provide, respectively, positive mass in a common shifted target and progressiveness. The latter uses Theorems 6.1 and 6.2. The uniformity estimates discussed below enter the proof of Theorem 6.2.\par
\section*{A false density assertion in the weak-mixing example}
Problem 1 (P01: Example 4.2, p. 11; Minor).  The displayed inclusion and the ensuing claim that the recurrence-time set has full density along every Følner sequence omit the distinguished zeroth coordinate of \(\sigma=\delta_a\otimes\mu^{\otimes k}\).\par
Equation reference: eq:mass. Put \(U_0=U_{k+1}=X\), as in the source. Direct evaluation gives \[\begin{aligned}
&\sigma\bigl((X\times U_1\times\cdots\times U_k)
 \cap T_\Delta^{-n}(U_1\times\cdots\times U_k\times X)\bigr) \notag\\
&\hspace{12mm}=
 \mathbf 1_{U_1}(T^n a)\,\mu(U_k)
 \prod_{i=1}^{k-1}\mu(U_i\cap T^{-n}U_{i+1}). \end{aligned}\] The condition \(\mu(U_0\cap T^{-n}U_1)>0\) on the source’s right-hand side is automatic when \(\mu(U_1)>0\); it cannot replace \(T^n a\in U_1\).\par
\subsection*{An admissible counterexample}
Take the two-sided fair Bernoulli shift \(X=\{0,1\}^{\mathbb Z}\), with product probability \(\mu\) and the left shift \(T\). Disjoint-coordinate cylinder functions are independent once their supports are sufficiently separated; approximation by cylinder functions proves mixing, hence weak mixing. Choose a point \(a\) generic for \(\mu\) along \(\Phi_N=\{1,\ldots,N\}\), and put \[k=2,\qquad U_1=U_2=\{x:x(0)=1\}.\] Such generic points exist by the ergodic theorem. For every \(n\geq1\), \[\mu(U_1\cap T^{-n}U_2)=\frac14,\qquad
 \mu(T^{-n}U_1)=\mu(U_2)=\frac12 .\] Consequently the right-hand intersection printed in Example 4.2 is all of \(\mathbb N\). By [eq:mass] (\#eq:mass), the left-hand positivity set is exactly \(\{n:a(n)=1\}\), since the measure there is \(\frac18\mathbf 1_{\{a(n)=1\}}\). That set has density \(1/2\), not one. This disproves the inclusion and the strong density claim under valid hypotheses. The same missing pointwise condition is visible already when \(k=1\).\par
\subsection*{Complete local repair and impact}
For general \(k\), write \[H=\{n:T^n a\in U_1\},\qquad
 D_i=\{n:\mu(U_i\cap T^{-n}U_{i+1})>0\},\quad 1\leq i<k.\] The positivity set is \(H\cap\bigcap_iD_i\). Weak mixing gives density one for each \(D_i\) along every Følner sequence. Along the particular sequence \(\Phi\) for which \(a\) is generic, the Portmanteau inequality for open \(U_1\) gives \[\liminf_N\frac{|H\cap\Phi_N|}{|\Phi_N|}\geq\mu(U_1)>0.\] Removing finitely many zero-density complements preserves this positive lower bound. Since \(|\Phi_N|\to\infty\), there are infinitely many admissible times. Thus \(\sigma\) is progressive, as intended. Moreover, \(\sigma(X\times(T^{-t}E)^k)=\mu(E)^k>0\) for every \(t\).\par
The correction restores the example’s required conclusion using its generic-point hypothesis. The false statement concerns an auxiliary full-density strengthening; it is not used as a replacement for the general proof of Theorem 4.9. The Minor classification follows from that verified dependency and local repair, not merely from the existence of some possible correction. Evidence is in P01, including the exact product calculation and the actually inspected p. 11 image.\par
\section*{Missing normalization in the quantitative estimate}
Problem 2 (P02: Theorem 6.6, p. 20, equation (6.6); Minor).  The estimate is stated for arbitrary bounded functions, while its constant depends only on \(k\) and its right-hand side is only homogeneous of degree one in the functions.\par
For \(k=3\), take a one-point probability system, with its identity transformation, and let \(\tau\) be the unique measure on its fourth power. This system is ergodic, \(\tau\) has the required invariance, and \(\tau_i=\mu\leq i\mu\) for \(i=1,2\). Set \(b(n)=1\) and \(f_1=f_2=M>0\). Every Følner average on the left of (6.6) equals \(M^2\). Every uniformity seminorm of a positive constant \(M\) equals \(M\), by the recursive definition. The claimed inequality is therefore \[M^2\leq C_3 M
 \qquad\hbox{for every }M>0,\] which is impossible for a constant \(C_3\) independent of \(M\). For \(k\geq3\), the same construction compares degrees \(k-1\) and one. This objection does not apply to the single-function case \(k=2\).\par
\subsection*{The normalized and homogeneous versions}
Equation reference: eq:homogeneous. The proof on pp. 24–25 invokes Theorem 6.9, whose functions are explicitly bounded by one. Supplying \(\|f_j\|_\infty\leq1\) is therefore a genuine missing hypothesis in the statement and that invocation. Once the normalized estimate is established, multilinearity gives the correct homogeneous form \[\begin{aligned}
&\limsup_N\left\|
 \frac 1{|\Phi_N|}\sum_{n\in\Phi_N}b(n)(1\otimes T^nf_1\otimes\cdots\otimes T^nf_{k-1}\otimes 1)
 \right\|_{L^2(\tau)} \notag\\
&\quad\leq C_k\|b\|_\infty
 \min_{1\leq i<k}
 \left(\|f_i\|_{U^k(\mu,T)}
       \prod_{\substack{1\leq j<k\\j\ne i}}\|f_j\|_\infty\right).
 \end{aligned}\] If all norms are positive, divide each function by its supremum norm, apply the unit-ball estimate, and multiply back. If one is zero, absolute continuity of the corresponding marginal with respect to \(\mu\), together with \(T\)-invariance of \(\mu\), makes every summand zero in \(L^2(\tau)\). Thus this endpoint also follows.\par
\subsection*{Dependency and severity}
Theorem 6.2 begins with \(\|f_i\|_\infty\leq1\) and uses this estimate to discard terms containing \(f_i-\mathbb E(f_i\mid Z_{k-1})\), whose supremum norm is at most two and whose \(U^k\) seminorm is zero. Dividing that difference by two, or applying [eq:homogeneous] (\#eq:homogeneous), gives exactly the required zero limit. It does not create an uncontrolled constant in a nonzero error estimate. The tensor argument of Lemma 6.10 and the normalized version of Theorem 6.9 supply the remaining bound, subject to the restricted-coefficient correction in Problem 3.\par
The printed uniform estimate is false. Its actual use in the core proof requires only the normalized zero-seminorm conclusion, and that use survives the demonstrated homogeneous correction. This specific boundary supports a local Minor classification; no failure of Theorem 6.2 is deduced from scaling. Evidence is in P02.\par
\section*{The multiple-average lemma needs coefficient restrictions}
Problem 3 (P03: Lemma 6.7(iii), p. 21; Minor).  The lemma permits arbitrary positive integers \(c_1,\ldots,c_k\) and all \(k\geq1\). Its asserted \(U^k\) control fails both for repeated coefficients and, independently, for the stated one-function endpoint.\par
\subsection*{Repeated coefficients}
Use the two-sided fair Bernoulli shift on \(\{-1,1\}^{\mathbb Z}\) and let \(f(x)=x(0)\). This system is ergodic and \(|f|=1\). For every \(h\geq1\), \[\int (T^h f)f\,d\mu=0.\] Ergodicity and the definition of \(U^2\) give \[\|f\|_{U^2}^4
 =\lim_{H\to\infty}\frac1H\sum_{h=1}^H
       \left|\int(T^hf)f\,d\mu\right|^2=0.\] Take \(k=2\), \(c_1=c_2=1\), and \(f_1=f_2=f\). Every term of the average is \((T^nf)^2=1\). Its \(L^2\) norm is one, whereas the proposed upper bound is zero, for every finite \(C\). All unit-ball hypotheses hold. The obstruction is a repeated linear form, not a missing numerical constant.\par
\subsection*{The one-function endpoint}
Take the two-point space \(\{0,1\}\) with uniform measure and \(T(x)=x+1\pmod2\). It is ergodic. For \(f(0)=1\), \(f(1)=-1\), \[\|f\|_{U^1(\mu,T)}=\left|\int f\,d\mu\right|=0.\] With \(k=1\) and \(c_1=2\), however, \(T^{2n}f=f\) for every \(n\). The average has \(L^2\) norm one. The correct one-function limit is \(\|\mathbb E(f\mid\mathcal I_{T^{c_1}})\|_2\), which need not be controlled by \(U^1(\mu,T)\).\par
\subsection*{Repair boundary and actual application}
For \(k\geq2\), the bound needed here is valid for pairwise distinct nonzero coefficients, with a limsup if existence of the norm limit is not separately invoked. Host’s Proposition 1 [host] bounds the average, for each distinguished \(i\), by the box seminorm with directions \(c_i\) and \(c_j-c_i\) for \(j\ne i\). All these directions are nonzero. The power comparison in Frantzikinakis–Kuca, Lemma 3.1 [fk], gives \[\|\cdot\|_{\text{box}(c_i,(c_j-c_i)_{j\ne i})}
 \leq R_i^{1/2^k}\|\cdot\|_{U^k(T)},\qquad
 R_i=|c_i|\prod_{j\ne i}|c_j-c_i|.\] The cited comparison is valid on general measure-preserving systems for degree at least two. Real/imaginary expansion, if needed, loses only a finite factor depending on \(k\). The interval estimate transfers to arbitrary one-dimensional Følner sets by decomposing them into consecutive runs: the proportion in runs of length below \(L\) is at most \(L\) times the boundary-to-volume ratio. First choose large \(L\), then let \(N\) increase. This proves the restricted bound used downstream, with a constant depending only on the fixed coefficients and \(k\). For \(k=1\), one may use \(c_1=1\), or retain the invariant-factor norm for \(T^{c_1}\); replacing it by the \(T\)-invariant norm without further hypotheses is invalid. The FK dependency checked here is its specified author manuscript, not an asserted copy of its journal final version.\par
In Theorem 6.9, pp. 22–23, the coefficients are \(c_i=c/i\), where \(c=\operatorname{lcm}(1,\ldots,k)\). They are positive, integral, and pairwise distinct. Its \(k=1\) case can be proved directly: \(\nu_1\leq\mu\) and equality of total masses imply \(\nu_1=\mu\), and the mean ergodic theorem gives the required \(U^1\) norm. Thus the actual application falls within the valid restricted statement. In turn, Theorem 6.6 uses Theorem 6.9 for \(k-1\) functions; its smallest case is covered by that direct one-function argument.\par
These counterexamples disprove the lemma as printed, including a unit-bounded example with zero right-hand side. The impact is confined here to overbroad auxiliary hypotheses because the specific downstream coefficients and endpoint have been checked. The Minor assessment does not rest on a rule that a supporting lemma must overturn the main theorem to be serious; it rests on the verified applicability of the corrected version to this proof. Evidence is in P03.\par
\section*{Positive checks and boundaries of the reconstruction}
\subsection*{Correspondence and positive-time genericity}
For Theorem 2.3, pp. 6–7, a two-sided symbolic shift with the positive coordinates of \(a\) equal to \(\mathbf 1_A\) gives the exact return set to the clopen cylinder \(E=\{x:x(0)=1\}\). Weak-star limits of orbit averages on long intervals of positive \(A\)-density are invariant: translating an interval changes only its endpoints. They give \(\mu(E)>0\) and are supported on the forward orbit closure. The ergodic-component reduction retains positive mass on \(E\).\par
The claimed genericity along some positive Følner sequence can also be seen directly. Choose an ergodic generic point \(y\) in that orbit closure. For longer and longer initial orbit segments of \(y\), continuity allows a forward translate of \(a\) to shadow the averages of finitely many continuous test functions. A diagonal choice over a countable dense family gives intervals in \(\mathbb N\) on which \(a\) is generic for that ergodic component. This validates the forward-time requirement without assuming that \(a\) is a typical initial-segment generic point.\par
\subsection*{Extraction from a progressive measure}
In Lemma 2.2, the recursive neighborhoods remain neighborhoods of their specified points; the common convergence subsequence meets every finite set of requirements. Peeling the selected indices in descending order places each distinct-term sum in its prescribed target. In Proposition 3.2, the actual p. 9 image contains the closure bars on the nested sets and their intersection. Their loss in extracted text is not an error in the paper. Using a max product metric gives product balls, so the rectangular definition of progressiveness applies. Positive-mass open intersections meet the support, and smaller balls with compactly nested closures give the limiting progression.\par
\subsection*{Disintegration, recurrence, and nilpotent structure}
The marginal estimate in Lemma 4.6 has a direct explanation. If \(\xi_i\) is \(T^i\)-invariant, then \(i^{-1}\sum_{j=0}^{i-1}T_*^j\xi_i\) is \(T\)-invariant and therefore equals the unique invariant measure of the ergodic pronilsystem. Positivity gives \(\xi_i\leq i m_{k-1}\). This also ensures that the integrals of conditional expectations in Definition 4.5 are independent of their almost-everywhere representatives. Disintegration then defines a positive probability measure \(\sigma\) with the stated marginals. Equivariance gives Lemma 4.7.\par
The use of a continuous topological factor includes intertwining on the relevant orbits, in particular at \(a\). Mere continuity of an almost-everywhere equivariant measurable map, without the topological factor property, would not justify evaluation at an exceptional point. The inspected lifting construction [sumsets] takes an orbit closure inside the product of the original system and its pronilfactors. The factor maps are coordinate projections and intertwine everywhere, including at the lifted point. The measurable graph map is used to identify the invariant measure and its ergodicity. Restricting the base to the orbit closure of \(a\) supplies the transitivity hypothesis. This validates the intended application without turning a purely almost-everywhere interpretation into an unproved finding.\par
In Theorem 6.3, setting \(c=\operatorname{lcm}(v_1,\ldots,v_k)\) and \(w_i=cu_i/v_i\) reduces the required rectangular intersection to a multiple-recurrence intersection for the measure-preserving map \(T_v\). For an arbitrary Følner sequence, residue classes have asymptotic proportion \(1/c\), and \(\Psi_N=\{n:cn\in\Phi_N\}\) is eventually nonempty and Følner. The factor \(1/c\) therefore preserves positivity. Furstenberg’s theorem applies without an ergodicity assumption on that system. The product-set hypothesis is essential; Example 6.4 correctly exhibits its failure for a nonrectangular subtorus.\par
\subsection*{The final approximation argument}
In Theorem 6.2, unit-bounded tensor products approximate a continuous \(G\) by finite linear combinations. A uniform approximation error \(\delta\) contributes at most \(2\delta\) to the difference of the two averages, without requiring \(\sigma\) to be \(T_\Delta\)-invariant. Conditional-expectation differences have zero \(U^k\) seminorm. The normalized forms discussed in Problems 2–3 consequently remove the unstructured factors. Marginal domination controls the continuous approximations to the remaining conditional expectations.\par
The pointwise-in-\(y\) estimate of Lemma 6.11 is particularly useful. For the quantity \(Z_N(y)\) in (6.24), it gives \(\limsup_N|Z_N(y)|\leq\varepsilon\) for every \(y\), and \(|Z_N(y)|\leq2\). Reverse Fatou applied to the squares yields \[\limsup_N\|Z_N\|_{L^2(\xi)}^2
 \leq\int\limsup_N |Z_N(y)|^2\,d\xi(y)
 \leq\varepsilon^2.\] This bypasses the p. 29 estimate that would otherwise need a factor two when bounding an \(L^2\) square by an \(L^1\) norm. The resulting error still tends to zero; no additional substantive problem is counted.\par
For (6.22), expand the square and use two copies of the nilpotent orbit space, with commuting actions corresponding to the two averaging variables and their common progression variable. Unique ergodicity of the joint orbit closure makes the joint Følner limit independent of order. For fixed progression parameter, the translation \(n\mapsto n-m\) changes a Følner average by a vanishing boundary error, so the inner difference tends to zero. For each fixed outer index, the needed inner limit exists: in the progression variable by orbit unique ergodicity, and in the averaging variable by the just-described vanishing. Uniform convergence when both indices are large then identifies both iterated limits, as in the bounded-Fubini principle of [bl]. This gives the required limit exchange and vanishing. The existence and structure of these pronilsystems remain explicit external inputs; their entire theory is not claimed re-proved here.\par
For \(k=1\), set \(\sigma=\delta_a\otimes\mu\) directly. Genericity sends the scalar average \(G(T^na)\) to \(\int G\,d\mu\), and the Følner mean ergodic theorem sends the other average to the same constant in \(L^2(\mu)\). This proves Theorem 6.2 at that endpoint without applying Theorem 6.6 outside its stated range.\par
\section*{Conclusion}
The reviewed version needs local correction of three mathematical statements. Problems P01–P03 are supported by exact counterexamples and specific repairs. Their uses in the main argument have been traced: no failure of the asserted density finite sums conclusion has been established within this scope. This is not a certification that no other issue exists.\par
Dimension | Recorded outcome\par
Report assessment | Minor revision within the inspected scope\par
Confirmed issues | Three local missing-condition or false-density statements\par
Core dependencies | Specific corrected uses, positive-time reductions, and bounded approximation/limit arguments checked\par
External limits | Full nilpotent structure theory and every cited proof not independently reconstructed\par
Scientific adjudication | Pending\par
Human reviews | None recorded\par
Author adoption | Unknown; no contact or adoption asserted\par
Independent acceptance | Not performed in this initial review\par
\section*{Supporting checks}
4ad90da06e2236c4e0db81bc0b7585ca3e159364ce39ebec29f4719a3ff13693.\par
The matching P01–P03 evidence manifests bind the original locations, proofs, counterevidence, repair limits, and actual source-page images. Any reviewer-authored exact calculation is supplementary to the analytic argument, not a finite-sample substitute for a universal proof. No author code was executed. Its absence is not a mathematical defect.\par
The final report is compiled from its exact current TeX using the previously authorized portable Tectonic runtime, in untrusted mode without shell escape. Real command lines, exit status, dependencies, logs, PDF and image hashes, and actual inspection of every final page are retained in the provenance. Prior drafts and operational failures are preserved.\par
Three existing native helpers supplied bounded checks; their task names retain 0107 because those existing tasks were reused, while their artifacts explicitly bind 0109. They are not independent acceptance reviewers. The external controller’s reported UI observation was Personal account, Chat, GPT-6, Pro. Neither label is asserted to be an independent measurement of the backend model. The earlier controller-contract helper was operational only and supplies no scientific evidence.\par
9 B. Kra, J. Moreira, F. K. Richter, and D. Robertson, The density finite sums theorem, exact acquired publisher PDF, 31 pp., DOI https://doi.org/10.1007/s00222-025-01371-8. The first-page header prints 2025; online date September 24, 2025. Springer Nature, publisher record for the same article, https://link.springer.com/article/10.1007/s00222-025-01371-8, consulted October 26, 2026; issue metadata 243 (2026), 1–31. B. Host, Ergodic seminorms for commuting transformations and applications, Studia Math. 195 (2009), 31–49, Proposition 1, printed p. 33; publisher PDF consulted. https://www.impan.pl/shop/en/publication/transaction/download/product/89663. N. Frantzikinakis and B. Kuca, Joint ergodicity for commuting transformations and applications to polynomial sequences, arXiv:2207.12288v2, December 18, 2024, Lemma 3.1, PDF pp. 20–21. https://arxiv.org/pdf/2207.12288v2. Only this dependency version is claimed inspected. B. Kra, J. Moreira, F. Richter, and D. Robertson, Infinite sumsets in sets with positive density, author manuscript dated July 20, 2022, local construction in Proposition 5.6 and Lemma 5.7. https://sites.math.northwestern.edu/\textasciitilde{}kra/papers/higher-sums.pdf. Published-numbered statements were separately located as Proposition 5.7 and Lemma 5.8; full version equivalence is not asserted. V. Bergelson and A. Leibman, Cubic averages and large intersections, author manuscript dated January 22, 2018, Lemma 1.1, p. 9. https://people.math.osu.edu/bergelson.1/BL\_Cubic.pdf. This consulted manuscript is not asserted identical to the 2015 chapter.\par
\end{document}
